\section{余下账本事件导航：western-jin}

\label{sec:remaining-event-anchors-western-jin}

本节为尚未在正文设置目标的账本记录建立直接跳转。锚点显示可核日期与事件名称；其解释仍应回到账本所列来源、置信提示及既有正文的区域和材料边界。

\subsection{事件导航与来源边界}

\eventanchor{JH-LEDGER-018-1}{晏平元年（306年）·“李雄在成都称帝，国号成，使巴蜀割据由起事集团转为帝国体制”的前期动员与资源准备}

\eventanchor{JH-LEDGER-018-2}{晏平元年（306年）·“李雄在成都称帝，国号成，使巴蜀割据由起事集团转为帝国体制”的命令、联盟或地方响应}

\eventanchor{JH-LEDGER-018-3}{晏平元年（306年）·李雄在成都称帝，国号成，使巴蜀割据由起事集团转为帝国体制}

\eventanchor{JH-LEDGER-018-4}{晏平元年（306年）·“李雄在成都称帝，国号成，使巴蜀割据由起事集团转为帝国体制”的执行中的空间与人员处置}

\eventanchor{JH-LEDGER-018-5}{晏平元年（306年）·“李雄在成都称帝，国号成，使巴蜀割据由起事集团转为帝国体制”的后续制度或军政调整}

\eventanchor{JH-LEDGER-018-6}{晏平元年（306年）·“李雄在成都称帝，国号成，使巴蜀割据由起事集团转为帝国体制”的异源材料与影响范围的复核}

\eventanchor{JH-LEDGER-020-1}{永凤元年（308年）·“刘渊称帝，正式以汉为国号并以平阳为政治中心”的前期动员与资源准备}

\eventanchor{JH-LEDGER-020-2}{永凤元年（308年）·“刘渊称帝，正式以汉为国号并以平阳为政治中心”的命令、联盟或地方响应}

\eventanchor{JH-LEDGER-020-3}{永凤元年（308年）·刘渊称帝，正式以汉为国号并以平阳为政治中心}

\eventanchor{JH-LEDGER-020-4}{永凤元年（308年）·“刘渊称帝，正式以汉为国号并以平阳为政治中心”的执行中的空间与人员处置}

\eventanchor{JH-LEDGER-020-5}{永凤元年（308年）·“刘渊称帝，正式以汉为国号并以平阳为政治中心”的后续制度或军政调整}

\eventanchor{JH-LEDGER-020-6}{永凤元年（308年）·“刘渊称帝，正式以汉为国号并以平阳为政治中心”的异源材料与影响范围的复核}

\eventanchor{JH-LEDGER-021-1}{永嘉三年（309年）·“刘聪、王弥等围攻洛阳，西晋京师依靠临时征兵和城防勉强守住”的前期动员与资源准备}

\eventanchor{JH-LEDGER-021-2}{永嘉三年（309年）·“刘聪、王弥等围攻洛阳，西晋京师依靠临时征兵和城防勉强守住”的命令、联盟或地方响应}

\eventanchor{JH-LEDGER-021-3}{永嘉三年（309年）·刘聪、王弥等围攻洛阳，西晋京师依靠临时征兵和城防勉强守住}

\eventanchor{JH-LEDGER-021-4}{永嘉三年（309年）·“刘聪、王弥等围攻洛阳，西晋京师依靠临时征兵和城防勉强守住”的执行中的空间与人员处置}

\eventanchor{JH-LEDGER-021-5}{永嘉三年（309年）·“刘聪、王弥等围攻洛阳，西晋京师依靠临时征兵和城防勉强守住”的后续制度或军政调整}

\eventanchor{JH-LEDGER-021-6}{永嘉三年（309年）·“刘聪、王弥等围攻洛阳，西晋京师依靠临时征兵和城防勉强守住”的异源材料与影响范围的复核}

\eventanchor{JH-LEDGER-022-1}{河瑞二年（310年）·“刘渊去世后刘聪经宫廷斗争继位，汉赵继续对西晋发动攻势”的前期动员与资源准备}

\eventanchor{JH-LEDGER-022-2}{河瑞二年（310年）·“刘渊去世后刘聪经宫廷斗争继位，汉赵继续对西晋发动攻势”的命令、联盟或地方响应}

\eventanchor{JH-LEDGER-022-3}{河瑞二年（310年）·刘渊去世后刘聪经宫廷斗争继位，汉赵继续对西晋发动攻势}

\eventanchor{JH-LEDGER-022-4}{河瑞二年（310年）·“刘渊去世后刘聪经宫廷斗争继位，汉赵继续对西晋发动攻势”的执行中的空间与人员处置}

\eventanchor{JH-LEDGER-022-5}{河瑞二年（310年）·“刘渊去世后刘聪经宫廷斗争继位，汉赵继续对西晋发动攻势”的后续制度或军政调整}

\eventanchor{JH-LEDGER-022-6}{河瑞二年（310年）·“刘渊去世后刘聪经宫廷斗争继位，汉赵继续对西晋发动攻势”的异源材料与影响范围的复核}

\eventanchor{JH-LEDGER-023-1}{永嘉四年（310年）·“司马越在东出避乱途中去世，其部众和辎重随即失去统一指挥”的前期动员与资源准备}

\eventanchor{JH-LEDGER-023-2}{永嘉四年（310年）·“司马越在东出避乱途中去世，其部众和辎重随即失去统一指挥”的命令、联盟或地方响应}

\eventanchor{JH-LEDGER-023-3}{永嘉四年（310年）·司马越在东出避乱途中去世，其部众和辎重随即失去统一指挥}

\eventanchor{JH-LEDGER-023-4}{永嘉四年（310年）·“司马越在东出避乱途中去世，其部众和辎重随即失去统一指挥”的执行中的空间与人员处置}

\eventanchor{JH-LEDGER-023-5}{永嘉四年（310年）·“司马越在东出避乱途中去世，其部众和辎重随即失去统一指挥”的后续制度或军政调整}

\eventanchor{JH-LEDGER-023-6}{永嘉四年（310年）·“司马越在东出避乱途中去世，其部众和辎重随即失去统一指挥”的异源材料与影响范围的复核}

\eventanchor{JH-LEDGER-024-1}{永嘉五年六月（311年）·“汉赵军攻陷洛阳，大量官民死散，西晋都城体系遭到毁灭性破坏”的前期动员与资源准备}

\eventanchor{JH-LEDGER-024-2}{永嘉五年六月（311年）·“汉赵军攻陷洛阳，大量官民死散，西晋都城体系遭到毁灭性破坏”的命令、联盟或地方响应}

\eventanchor{JH-LEDGER-024-3}{永嘉五年六月（311年）·汉赵军攻陷洛阳，大量官民死散，西晋都城体系遭到毁灭性破坏}

\eventanchor{JH-LEDGER-024-4}{永嘉五年六月（311年）·“汉赵军攻陷洛阳，大量官民死散，西晋都城体系遭到毁灭性破坏”的执行中的空间与人员处置}

\eventanchor{JH-LEDGER-024-5}{永嘉五年六月（311年）·“汉赵军攻陷洛阳，大量官民死散，西晋都城体系遭到毁灭性破坏”的后续制度或军政调整}

\eventanchor{JH-LEDGER-024-6}{永嘉五年六月（311年）·“汉赵军攻陷洛阳，大量官民死散，西晋都城体系遭到毁灭性破坏”的异源材料与影响范围的复核}

\eventanchor{JH-LEDGER-025-1}{永嘉五年（311年）·“怀帝在洛阳陷落时被汉赵俘获并押往平阳，西晋皇统陷入屈辱性危机”的前期动员与资源准备}

\eventanchor{JH-LEDGER-025-2}{永嘉五年（311年）·“怀帝在洛阳陷落时被汉赵俘获并押往平阳，西晋皇统陷入屈辱性危机”的命令、联盟或地方响应}

\eventanchor{JH-LEDGER-025-3}{永嘉五年（311年）·怀帝在洛阳陷落时被汉赵俘获并押往平阳，西晋皇统陷入屈辱性危机}

\eventanchor{JH-LEDGER-025-4}{永嘉五年（311年）·“怀帝在洛阳陷落时被汉赵俘获并押往平阳，西晋皇统陷入屈辱性危机”的执行中的空间与人员处置}

\eventanchor{JH-LEDGER-025-5}{永嘉五年（311年）·“怀帝在洛阳陷落时被汉赵俘获并押往平阳，西晋皇统陷入屈辱性危机”的后续制度或军政调整}

\eventanchor{JH-LEDGER-025-6}{永嘉五年（311年）·“怀帝在洛阳陷落时被汉赵俘获并押往平阳，西晋皇统陷入屈辱性危机”的异源材料与影响范围的复核}

\eventanchor{JH-LEDGER-026-1}{建兴元年（313年）·“汉赵杀害被俘的怀帝，长安的司马邺被拥立为愍帝”的前期动员与资源准备}

\eventanchor{JH-LEDGER-026-2}{建兴元年（313年）·“汉赵杀害被俘的怀帝，长安的司马邺被拥立为愍帝”的命令、联盟或地方响应}

\eventanchor{JH-LEDGER-026-3}{建兴元年（313年）·汉赵杀害被俘的怀帝，长安的司马邺被拥立为愍帝}

\eventanchor{JH-LEDGER-026-4}{建兴元年（313年）·“汉赵杀害被俘的怀帝，长安的司马邺被拥立为愍帝”的执行中的空间与人员处置}

\eventanchor{JH-LEDGER-026-5}{建兴元年（313年）·“汉赵杀害被俘的怀帝，长安的司马邺被拥立为愍帝”的后续制度或军政调整}

\eventanchor{JH-LEDGER-026-6}{建兴元年（313年）·“汉赵杀害被俘的怀帝，长安的司马邺被拥立为愍帝”的异源材料与影响范围的复核}

\eventanchor{JH-LEDGER-029-1}{建兴四年十一月（316年）·“刘曜攻陷长安，愍帝出降，西晋在北方的最后朝廷覆灭”的前期动员与资源准备}

\eventanchor{JH-LEDGER-029-2}{建兴四年十一月（316年）·“刘曜攻陷长安，愍帝出降，西晋在北方的最后朝廷覆灭”的命令、联盟或地方响应}

\eventanchor{JH-LEDGER-029-3}{建兴四年十一月（316年）·刘曜攻陷长安，愍帝出降，西晋在北方的最后朝廷覆灭}

\eventanchor{JH-LEDGER-029-4}{建兴四年十一月（316年）·“刘曜攻陷长安，愍帝出降，西晋在北方的最后朝廷覆灭”的执行中的空间与人员处置}

\eventanchor{JH-LEDGER-029-5}{建兴四年十一月（316年）·“刘曜攻陷长安，愍帝出降，西晋在北方的最后朝廷覆灭”的后续制度或军政调整}

\eventanchor{JH-LEDGER-029-6}{建兴四年十一月（316年）·“刘曜攻陷长安，愍帝出降，西晋在北方的最后朝廷覆灭”的异源材料与影响范围的复核}
